\section{Quenched partition readout and boundary arms}
\label{sec:readout}

The marked estimates concern random functions in a partition integrand.
We now recover the spin correlation normalized separately in each
random environment.  Two features of this passage matter.  The charge of
an entire source list need not depend on local randomness, so independence
will be used only after replacing it by the charge of the principal
slot.  Also, large terminal background interactions cannot be discarded
globally.  We remove only the interactions meeting a growing buffer
around the two marks and use a conditional reference estimate inside
that buffer.

Throughout this section $D>0$ is fixed and sufficiently small, and
$t=t_*(D)$ is the tuned parameter of Proposition~\ref{flow:tuned}.
The expectation $\E$ includes the independent decision seeds of the
blocking construction.  All assertions about growth as the scale index
tends to infinity are uniform in the locations of the marks.  Constants
may depend on $D$ and on the fixed geometric parameters.  In particular,
no uniform large-distance asymptotic as $D\downarrow0$ is asserted.

\subsection{The exact identity and the two local amplitudes}

Let $y_1,y_2\in\Z^2$ and put
$\ell=\lVert y_1-y_2\rVert_\infty$.  Fix $F$ larger than all principal
support and dependence halos and all contact distances in one blocking
step.  Increase it, if necessary, so that the same separation holds in
output-scale units for every step considered below.  For sufficiently
large $\ell$, let $J$ be the largest integer such that
\begin{equation}
 F\le \ell/r_J<LF.
 \label{read:stop}
\end{equation}
Thus $J=\log_L\ell+O(1)$.  None of our estimates requires the ratio
$\ell/r_J$ to converge, or either mark to occupy a prescribed position
inside its grid cell.

Use two commuting square-free variables $\zeta_1,\zeta_2$, so
$\zeta_1^2=\zeta_2^2=0$.  Start the log integrand with
$\zeta_1\sigma_{y_1}+\zeta_2\sigma_{y_2}$; its initial joint coefficient
is zero.  At scale $j$ write the source-dependent part of the transformed
log integrand and its saved scalar as
\begin{equation}
 r_j^{-d}\sum_{a=1}^2\zeta_a\bigl(U_{j,a}+k_{j,a}\bigr)
 +r_j^{-2d}\zeta_1\zeta_2\bigl(U_{j,12}+k_{j,12}\bigr),
 \qquad d=\tfrac18.
 \label{read:jet}
\end{equation}
Here $U_{j,a}=\sum_i u_{j,a,i}$ is the sum of the scaled singleton
list, $U_{j,12}$ is the sum of the scaled joint list, and $k_{j,a},k_{j,12}$
are the accumulated scalar coefficients with the indicated scaling.
Saved scalars retain their carriers until all estimates involving those
carriers have been made.  The individual singleton lists coincide with
those obtained by running one source alone: set the other square-free
variable equal to zero and then differentiate in the chosen variable.

We apply the exact blocking identity first in a free box
$\Lambda_N$ containing the marks and all relevant buffers.  Write
$\langle\cdot\rangle_{J,N,\omega}$ for expectation in its actual
terminal background: the reference measure multiplied by the
transformed positive weight and divided by its own source-zero
partition integral.  The saved source-zero scalar cancels in this
ratio, however large it is.  Use the plane centering convention also
for source coefficients near the free boundary.  Spin flip then makes
$k_{J,1}=k_{J,2}=0$.  Expanding the exponential of
\eqref{read:jet} therefore gives the exact identity
\begin{equation}
 r_J^{2d}\langle\sigma_{y_1}\sigma_{y_2}\rangle^{\mathrm{free}}_{N,\omega}
 =\langle U_{J,1}U_{J,2}+U_{J,12}\rangle_{J,N,\omega}+k_{J,12}.
 \label{read:exact-normalized}
\end{equation}
This is an identity at each environment and each choice of decision
seeds.  In particular it computes the quenched correlation and not a
ratio of disorder-averaged partition integrals.  For every function
$Q$, the bound
$|\langle Q\rangle_{J,N,\omega}|\le\norm Q_\infty$ is independent
of the terminal interactions outside its support.  This is the reason
that the sup-norm errors below remain usable after readout.

We next estimate the plane lists.  Their use in
\eqref{read:exact-normalized} will be justified by the fixed-step
finite-volume passage below.
Let $u_{j,y,\mathrm{pr}}$ denote the principal slot of the singleton
list at $y$.  In addition to its total normalized charge $\alpha_{j,y}$
from Section~\ref{sec:marked}, define
\begin{equation}
 \widetilde\alpha_{j,y}
       =r_j^{-d}\ell_s(u_{j,y,\mathrm{pr}}),
 \qquad
 w_{j,y}=r_j^d P_{Hr_j,c_j(y)}\sigma_y.
 \label{read:local-charge}
\end{equation}
The line representative $w_{j,y}$ occupies the principal slot, has
normalized charge one, and has uniformly bounded sup norm.  The charge
bound of Lemma~\ref{ref:charges}, with exponential weights paying
all carrier-size factors, shows that the charge functional is bounded
on a singleton list.  The nonprincipal slots of
$u_{j,y}-\alpha_{j,y}w_{j,y}$ are exactly those of $u_{j,y}$.
Consequently Proposition~\ref{mark:kernel} gives, for each of the fixed
even powers $q$ used below and $j\ge1$,
\begin{align}
 \norm{\alpha_{j,y}-\widetilde\alpha_{j,y}}_{L^q}
   &\le C A_{j,q}(y)^{1/q}s_j,\\
 \norm{u_{j,y}-\widetilde\alpha_{j,y}w_{j,y}}_{v,q}
   &\le C A_{j,q}(y)^{1/q}s_j,
 \label{read:principal-replacement}
\end{align}
where $A_{j,q}(y)=\E\alpha_{j,y}^q$ and the weighted norm includes
slotwise function sup norms.  Summation over slots preserves these
bounds because the weight has the counting slack of
Section~\ref{sec:blocking}.

The random variable $\widetilde\alpha_{J,y}$ is measurable with respect
to the environment and all decision seeds in the fixed principal
footprint at $y$.  Those footprints at $y_1,y_2$ are disjoint by the
choice of $F$, including the safety layer for edge variables.  Hence
$\widetilde\alpha_{J,y_1}$ and $\widetilde\alpha_{J,y_2}$ are exactly
independent.  This assertion concerns the two amplitudes; it makes no
independence assertion about the Gibbs state in which they will be read.
Propositions~\ref{mark:moments} and~\ref{mark:first}, together with
\eqref{read:principal-replacement}, imply, for large $J$,
\begin{equation}
 \begin{split}
  c\le \E\widetilde\alpha_{J,y}\le C,\\
  \norm{\widetilde\alpha_{J,y}}_{L^2}
       =J^{1/16+o(1)},\\
  \norm{\widetilde\alpha_{J,y}}_{L^4}
       =J^{3/16+o(1)}.
 \end{split}
 \label{read:local-moments}
\end{equation}
For the first assertion the error is at most
$C A_{J,2}^{1/2}s_J=o(1)$; for the other two, the error in $L^q$ is
$o(A_{J,q}^{1/q})$.  Thus, with
$Z_J=\widetilde\alpha_{J,y_1}\widetilde\alpha_{J,y_2}$,
\begin{equation}
 c\le\E Z_J\le C,
 \qquad \E|Z_J|\le J^{1/8+o(1)},
 \qquad \norm{Z_J}_{L^2}=J^{1/8+o(1)},
 \qquad \norm{Z_J}_{L^4}=J^{3/8+o(1)}.
 \label{read:product-moments}
\end{equation}
The mean factors even if the amplitudes are signed.  The upper bound
for $\E|Z_J|$ also follows from the $L^2$ bound.

\begin{lemma}[Separated joint-source terms]
\label{read:merge}
At the stopping scale \eqref{read:stop}, the sum of the sup norms of the
joint-source list and the absolute values of its saved scalar
increments, including their subsequent scaling, has $L^2$ norm at most
\begin{equation}
 C s_J J^{3/8+o(1)}=J^{-1/8+o(1)}.
 \label{read:merge-bound}
\end{equation}
Replacing $U_{J,1}U_{J,2}$ by $Z_Jw_{J,y_1}w_{J,y_2}$ has the same
$L^2$ bound on its sup norm.
\end{lemma}

\begin{proof}
A newly generated coefficient containing both labels is bilinear in
the two singleton lists.  Apply the connected coefficient bounds of
Lemma~\ref{block:tree}, retaining the output carrier containing
both marks.  At zero background degree a pair of principal slots does
not merge: their input carriers, averaging squares, and contact
neighborhoods are disjoint, also in output units, by our choice of $F$.
Their transformed moment therefore factors and contributes zero to
the joint logarithmic coefficient.  At the first step the two input
sources are principal.  At later steps a zero-background merge must
therefore contain a remainder from
\eqref{read:principal-replacement}; H\"older's inequality in $L^4$
gives a factor $s_j$ in its $L^2$ cost.

For ordinary analytic terms of positive background degree, retain one
background factor before bounding the others by the small thresholds
and the deterministic tail bounds.  To spell out why the retained
factor is small despite dependence, suppress the slot weights
and write $X_a=\norm{u_{j,a,i_a}}$ and $H=\norm{h_m}$.
The mixed bound of Proposition~\ref{norm:source}, for source power four,
and H\"older interpolation with respect to the measure $X_a^4\,\dd\Prob$
give
\[
 \E(X_a^4H^2)
 \le \bigl(\E X_a^4H^p\bigr)^{2/p}
       \bigl(\E X_a^4\bigr)^{1-2/p}
 \le C U_a^4s_j^2,
 \qquad U_a=C A_{j,4}(y_a)^{1/4}.
\]
It follows that
\[
 \norm{X_1X_2H}_{L^2}^2
 \le\bigl(\E X_1^4H^2\bigr)^{1/2}
      \bigl(\E X_2^4H^2\bigr)^{1/2}
 \le C U_1^2U_2^2s_j^2.
\]
The slot weights interpolate in the same way.

There is a separate estimate for changes caused by a large-field
decision.  The actual map and its ordinary analytic comparison agree
unless one of the finitely many local badness tests in that term
succeeds.  For a tested principal slot $m$, let
$B_m=\{\norm{h_m}>\delta\}$; a random threshold in
$[\delta,2\delta]$ can only reduce this event.  The same mixed bound
gives
\[
 \E(X_a^4\one_{B_m})
 \le\delta^{-p}\E(X_a^4\norm{h_m}^p)
 \le C U_a^4s_j^p,
 \qquad
 \norm{X_1X_2\one_{B_m}}_{L^2}
 \le C U_1U_2s_j^{p/2}.
\]
Normalized tilted moments, and their joint logarithmic derivatives,
are bounded by a fixed multiple of $X_1X_2$; their bounds contain no
exponential of a large principal norm.  Thus the last display also
bounds the decision-dependent terms.  The number and positions of
tested centers cost only the polynomial carrier factors already paid
by the tree weights.  Summing the tests and the analytic coefficients
as in Lemma~\ref{block:tree} gives a fresh joint coefficient of
cost at most $C U_1U_2s_j$, since $p/2>1$.  All additional small
factors and placements are paid by the thresholds and exponential
slack.  This separates the smallness supplied by a retained analytic
factor from the smallness supplied by the rare decision event.

An animal containing both marks has length at least
$c\ell/r_{j+1}$.  Retain one further fixed positive exponential weight
on the output before summing it or saving a scalar.  The fresh cost is
therefore at most
\begin{equation}
 C A_{j,4}(y_1)^{1/4}A_{j,4}(y_2)^{1/4}s_j
       \exp(-c\ell/r_{j+1}).
 \label{read:fresh-merge}
\end{equation}
The finitely many initial-scale adjustments obey this bound with a
fixed constant.  Once generated, a joint coefficient is transported
linearly in that coefficient.  The deterministic source bound, with
normalized tilts and the fixed source rescaling, bounds its total
weighted sup cost by a constant $K$ per step, uniformly in the sizes
of the principal background fields.  Scalars obey the same bound
when their rescaling is included.  Thus their accumulated cost is
bounded by the sum of \eqref{read:fresh-merge} multiplied by
$K^{J-j-1}$, over $0\le j<J$.

For completeness, this sum is uniform over the whole stopping window.
For $j\ge J/2$, $s_j\le C_Ds_J$ and
$A_{j,4}(y_a)^{1/4}\le C_{D,\eta}J^{3/16+\eta}$ for every $\eta>0$.
Writing $m=J-j-1$, the remaining numerical series is bounded by
\[
 \sum_{m\ge0}K^m\exp(-cF L^m)<\infty.
\]
For $j<J/2$, the same exponential contains a factor at most
$\exp(-cF L^{J/2-1})$, which pays $K^J$ and any polynomial bound
on the earlier marked moments.  This proves
\eqref{read:merge-bound}, with arbitrarily small fixed power losses
in place of its $o(1)$.

Finally set $R_a=U_{J,a}-\widetilde\alpha_{J,y_a}w_{J,y_a}$.
The identity
\[
 U_{J,1}U_{J,2}-Z_Jw_{J,y_1}w_{J,y_2}
 =R_1U_{J,2}+
   \widetilde\alpha_{J,y_1}w_{J,y_1}R_2
\]
and \eqref{read:principal-replacement} give the same $L^2$ bound by
H\"older with fourth moments.  No independence is used in this last
estimate.
\end{proof}

\subsection{The fixed-step passage from boxes to plane lists}

We explain the finite-volume passage before using the stochastic
plane bounds in \eqref{read:exact-normalized}.  Fix the marks and the
number $J$ of steps.  Proposition~\ref{block:finite} bounds the
principal functions deterministically at each such step, uniformly
in $N$.  For the fixed source labels, Lemma~\ref{block:tree} gives
bounded weighted operator costs for transport and for each merge,
also under the normalized bad-field averages.  Starting from the
bounded microscopic sources, iteration through these finitely many
steps therefore bounds the source lists and their saved marked
increments deterministically, uniformly in $N$; scalar carriers are
retained when applying the weights.  With any of the fixed exponential
weights, the tails of the
source lists and of background lists touching a fixed region are
summable, with constants allowed to depend on $J$.  Given a carrier
cutoff, all remaining footprints and all their decision neighborhoods
lie in a fixed finite set.  Once $N$ contains that set, their
coefficients equal the plane coefficients exactly under the common
coupling of environment and seeds.  First use this equality for the
finitely many retained slots and increments, and then let the carrier
cutoff increase.  The deterministic tail bounds show that the
corresponding errors tend to zero in every fixed $L^q$ used here.
This argument also applies to the carrier-weighted marked scalar
increments before they are summed.  Odd local increments vanish by
spin flip; in the free box the parity statement holds for the entire
construction.

In this order of limits, Lemma~\ref{read:merge} and
\eqref{read:principal-replacement} imply
\begin{equation}
 r_J^{2d}\langle\sigma_{y_1}\sigma_{y_2}\rangle^{\mathrm{free}}_{N,\omega}
 =Z_J\langle w_{J,y_1}w_{J,y_2}\rangle_{J,N,\omega}
       +\mathcal R_{J,N},
 \quad
 \limsup_{N\to\infty}\norm{\mathcal R_{J,N}}_{L^2}
       \le C s_JJ^{3/8+o(1)}.
 \label{read:amplitude-readout}
\end{equation}
Here replacement of the finite-volume principal charges by the plane
charges has zero error for all sufficiently large $N$.  We do not need
the terminal measures themselves to converge: a uniform estimate of
the expectation in \eqref{read:amplitude-readout} will suffice.

\subsection{Removing a buffer and taking the two disorder moments}

The functions $w_{J,y_1},w_{J,y_2}$ are supported in a region of bounded
diameter in $r_J$ units.  Enclose it by a square buffer with margin
\begin{equation}
 b=J^{4/25}
 \label{read:buffer-radius}
\end{equation}
terminal cells.  Figure~\ref{read:geometry} shows the principal random
footprints and the surrounding region used to compare the spin
expectation with its reference value.
\begin{figure}[tbp]
\centering
\begin{tikzpicture}[x=1cm,y=1cm,font=\small,
  contour/.style={draw=blue!55!black,line width=.8pt},
  footprint/.style={draw=blue!55!black,densely dashed,line width=.7pt},
  guide/.style={draw=black!45,densely dotted},
  dimension/.style={draw=black!65,{Stealth[length=1.6mm]}-{Stealth[length=1.6mm]}}]
  \fill[black!3] (-3.15,-3.15) rectangle (3.15,3.15);
  \fill[white] (-2.60,-1.5) rectangle (2.60,1.5);
  \draw[black!65,line width=.8pt] (-3.15,-3.15) rectangle (3.15,3.15);
  \node[anchor=south] at (0,3.28) {Terminal buffer};
  \draw[guide] (-2.60,-1.5) rectangle (2.60,1.5);

  \foreach \cx in {-1.5,1.5} {
    \draw[footprint,fill=blue!3] (\cx-1,-1) rectangle (\cx+1,1);
    \draw[contour] (\cx-.45,-.45) rectangle (\cx+.45,.45);
  }
  \fill (-1.41,.05) circle[radius=1.5pt];
  \fill (1.59,.05) circle[radius=1.5pt];
  \node[anchor=north] at (-1.41,-.03) {$y_1$};
  \node[anchor=north] at (1.59,-.03) {$y_2$};
  \node[anchor=south] at (-1.5,.50) {$w_{J,y_1}$};
  \node[anchor=south] at (1.5,.50) {$w_{J,y_2}$};
  \node[anchor=north] at (-1.5,-1.05) {$\widetilde\alpha_{J,y_1}$};
  \node[anchor=north] at (1.5,-1.05) {$\widetilde\alpha_{J,y_2}$};

  \draw[dimension] (2.80,1.5) -- (2.80,3.15)
    node[midway,left,fill=black!3,inner sep=2pt] {$b r_J$};
  \draw[guide] (2.60,1.5) -- (2.97,1.5);
  \draw[guide] (-1.41,-1.55) -- (-1.41,-2.08);
  \draw[guide] (1.59,-1.55) -- (1.59,-2.08);
  \draw[dimension] (-1.41,-1.92) -- (1.59,-1.92)
    node[midway,below,fill=black!3,inner sep=2pt] {$\ell$};
  \node at (0,-2.65) {$F\le \ell/r_J<LF$};

  \draw[contour] (3.60,1.62) -- (4.12,1.62);
  \node[anchor=west,align=left,text width=3.15cm] at (4.25,1.62)
    {Reference contours\\supporting $w_{J,y_a}$};
  \draw[footprint] (3.60,.35) -- (4.12,.35);
  \node[anchor=west,align=left,text width=3.15cm] at (4.25,.35)
    {Random footprints\\of the principal slots};
  \draw[guide] (3.60,-.95) -- (4.12,-.95);
  \node[anchor=west,align=left,text width=3.15cm] at (4.25,-.95)
    {A bounded insertion\\enclosure};
  \node[anchor=west,align=left,text width=3.7cm,font=\footnotesize]
    at (3.60,-2.45) {Schematic; not to scale.};
\end{tikzpicture}
\caption{Geometry of the terminal readout.  The solid inner squares are
the reference contours defining $w_{J,y_a}$, centered at the cells
$c_J(y_a)$.  The dashed principal footprints record the disorder and
decision seeds determining $\widetilde\alpha_{J,y_a}$; their disjointness,
including the safety layer for edge variables, gives independence of the
two amplitudes.  The dotted enclosure has bounded diameter in terminal
cell units.  A surrounding square provides margin of order $b r_J$,
where $b=J^{4/25}$.  Removing terminal interactions whose spin supports
meet this square leaves a reference conditional expectation inside it.
The two spin observables are still evaluated jointly in that
expectation.  Shapes and relative lengths are schematic.}
\label{read:geometry}
\end{figure}


Let $T_{J,N}$ be the sum of the sup norms of all
terminal background slots whose spin supports meet the buffer or its
interior.  We may bound it by the larger sum over slots whose recorded
carriers meet that region.  There are $O(b^2)$ possible cells at which
to root such a slot.  The exponential weights pay the choices of
animals and any distance of their anchors from those cells.  The
background $L^p$ bound from Proposition~\ref{flow:tuned} and Minkowski's
inequality consequently give
\begin{equation}
 \limsup_{N\to\infty}\norm{T_{J,N}}_{L^p}
    \le C b^2s_J=O_D(J^{-9/50}).
 \label{read:buffer-cost}
\end{equation}
Finite-box errors vanish by the preceding fixed-step argument.  We
have chosen the fixed background moment $p$ sufficiently large; any
$p>25$ suffices for the probability estimates below, and the earlier
construction allows much larger values.

Put $Q_J=w_{J,y_1}w_{J,y_2}$ and
$q_J=\langle Q_J\rangle_0$, where the latter is plane reference
expectation.  If $T_{J,N}\le1$, delete these background slots.
For a bounded observable and a perturbation $W$ with
$\norm W_\infty\le T$, differentiation of normalized expectations
along $\exp(tW)$, $0\le t\le1$, gives
\[
 \left|\frac{\dd}{\dd t}\langle Q\rangle_t\right|
   =|\Cov_t(Q,W)|\le2\norm Q_\infty T.
\]
Thus deletion changes the expectation of $Q_J$ by at most $CT_{J,N}$.
Conditionally on spins outside the buffer, the remaining law inside
is the reference Ising law with some boundary data.  The product
$Q_J$ is even and has bounded sup norm.  Its support diameter in
terminal units is bounded uniformly in \eqref{read:stop}.  The even
conditional reference estimate of Proposition~\ref{ref:operator}
therefore changes its expectation from the plane value by at most
$C/b$, uniformly in those outer data.  On the complementary event we
use the uniform bound $C$ on both expectations.  In particular, if
\[
 B_{J,N}=\{T_{J,N}>J^{-7/50}\},
 \qquad
 \Delta_{J,N}=\langle Q_J\rangle_{J,N,\omega}-q_J,
\]
then
\begin{equation}
 |\Delta_{J,N}|\le C(J^{-7/50}+J^{-4/25})
     \quad\hbox{on }B_{J,N}^{\mathrm c},
 \qquad |\Delta_{J,N}|\le C,
 \qquad
 \limsup_{N\to\infty}\Prob(B_{J,N})\le C_DJ^{-p/25}.
 \label{read:buffer-good}
\end{equation}
The probability bound follows from Markov's inequality and
\eqref{read:buffer-cost}; it requires no independence between
$T_{J,N}$ and the amplitudes.

Here are the distinct estimates needed for the first and second
moments.  On $B_{J,N}^{\mathrm c}$, \eqref{read:product-moments} gives
\begin{align*}
 \E\bigl[|Z_J\Delta_{J,N}|\one_{B_{J,N}^{\mathrm c}}\bigr]
     &\le J^{1/8-7/50+o(1)}=J^{-3/200+o(1)},\\
 \norm{Z_J\Delta_{J,N}\one_{B_{J,N}^{\mathrm c}}}_{L^2}
     &\le J^{1/8-7/50+o(1)}.
\end{align*}
On the exceptional event, H\"older's inequality gives, after the
finite-volume limit,
\begin{align*}
 \limsup_{N\to\infty}
 \E\bigl[|Z_J\Delta_{J,N}|\one_{B_{J,N}}\bigr]
     &\le C\norm{Z_J}_{L^2}J^{-p/50}
       =J^{1/8-p/50+o(1)},\\
 \limsup_{N\to\infty}
 \norm{Z_J\Delta_{J,N}\one_{B_{J,N}}}_{L^2}
     &\le C\norm{Z_J}_{L^4}J^{-p/100}
       =J^{3/8-p/100+o(1)}.
\end{align*}
The first error tends to zero.  For some fixed $\eta>0$ the second
is $O(J^{1/8-\eta+o(1)})$, since $p>25$.  Together with
\eqref{read:amplitude-readout}, these estimates prove
\begin{equation}
 \begin{split}
 \limsup_{N\to\infty}
 \E\left|r_J^{2d}
  \langle\sigma_{y_1}\sigma_{y_2}\rangle^{\mathrm{free}}_{N,\omega}
       -q_JZ_J\right|&=o(1),\\
 \limsup_{N\to\infty}
 \norm{r_J^{2d}
  \langle\sigma_{y_1}\sigma_{y_2}\rangle^{\mathrm{free}}_{N,\omega}
       -q_JZ_J}_{L^2}
       &=O(J^{1/8-\eta+o(1)}).
 \end{split}
 \label{read:two-errors}
\end{equation}
Decrease $\eta$ if necessary to include the joint-source error
$J^{-1/8+o(1)}$.  The separate $L^1$ estimate is useful: a relative
$L^2$ estimate alone would not recover the bounded first moment.

The two reference averaging squares defining $w_{J,y_a}$ are
disjoint.  Conditional independence in the reference model, given
their contours, and the tower property yield the exact equality
\begin{equation}
 q_J=r_J^{2d}\langle\sigma_{y_1}\sigma_{y_2}\rangle_0.
 \label{read:reference-pair}
\end{equation}
The separated reference correlation limit, uniformly over the
compact window \eqref{read:stop} and over directions, gives
$0<c\le q_J\le C$.  Thus arbitrary integer separations and grid-cell
positions are covered.

\begin{proposition}[Disorder moments at tuning]
\label{read:moments}
Let $C_{*,\omega}(y_1,y_2)$ be the free infinite-volume spin correlation
at $t_*(D)$, and let $J$ be determined by \eqref{read:stop}.  For every
fixed sufficiently small $D>0$, uniformly in the two marks as their
separation tends to infinity,
\begin{equation}
 r_J^{2d}\E C_{*,\omega}(y_1,y_2)\asymp1,
 \qquad
 r_J^{4d}\E C_{*,\omega}(y_1,y_2)^2
       =J^{1/4+o(1)}.
 \label{read:moment-conclusion}
\end{equation}
\end{proposition}

\begin{proof}
For each fixed separation, free ferromagnetic correlations increase
as the volume increases and remain in $[0,1]$.  Their limit exists,
and bounded convergence passes both disorder moments to the limit.
Here $J$ and $r_J$ are fixed during this passage.  The first line of
\eqref{read:two-errors}, \eqref{read:product-moments}, and
\eqref{read:reference-pair} give the first assertion, including a
positive lower bound.  For the second, the triangle inequality and
its reverse in $L^2$ give
\[
 \left|\norm{r_J^{2d}C_{*,\omega}(y_1,y_2)}_{L^2}
             -q_J\norm{Z_J}_{L^2}\right|
 \le O(J^{1/8-\eta+o(1)}).
\]
The main term is $J^{1/8+o(1)}$, so squaring proves the second
assertion.  All finite-volume identities hold for each choice of
seeds, whereas the original correlation depends only on the bonds.
Averaging over the seeds thus leaves the asserted disorder moments
unchanged.
\end{proof}

\subsection{A finite-box estimate with wired boundary}

The identification of the transition also needs an estimate in boxes
whose boundary is fixed to plus.  Such a boundary can create an odd
saved scalar, so it is important to keep the carrier of that scalar
until its distance from the central mark has been used.

\begin{proposition}[Wired arm at tuning]
\label{read:arm}
For any square box of radius $n$ centered at a lattice site $y$, let
$\phi^w_{n,*,\omega}$ be the FK--Ising law at $t_*(D)$ with wired
boundary.  Then
\begin{equation}
 \E\left[\phi^w_{n,*,\omega}
    (y\longleftrightarrow\partial\Lambda_n(y))^{400}\right]
       \le n^{-400/8+o(1)}.
 \label{read:arm-bound}
\end{equation}
The estimate is uniform in $y$.  It also holds if the fixed wall is
one lattice layer inside or outside the indicated square, including
the convention with bonds crossing to an exterior wired layer.
\end{proposition}

\begin{proof}
Under the Edwards--Sokal coupling~\cite{EdwardsSokal1988}, the wired
arm probability equals the spin expectation with the corresponding
plus wall.  Start a single source at $y$ and let $J_n$ be the largest
integer satisfying
\begin{equation}
 r_{J_n}\le n/(\log n)^2.
 \label{read:arm-stop}
\end{equation}
Thus $r_{J_n}\asymp n/(\log n)^2$, $J_n=O(\log n)$, and the principal
halo remains far from the wall for every $j\le J_n$.

Use the finite-box blocking rule, which forbids cuts near the wall
and leaves those interactions in the lists.  Every singleton term
whose full footprint stays in the bulk equals its plane counterpart.
By the weighted $L^{400}$ source estimate and
Proposition~\ref{mark:moments}, the sum of the sup costs of these
bulk terms at the final scale has $L^{400}$ norm bounded by a fixed
power of $1+J_n$.  Indeed
$A_{j,400}^{1/400}=j^{399/16+o(1)}$, and summing the exponential
slot weights only changes its constant.  A local odd scalar is zero:
before the boundary spins are evaluated, the bulk kernels and the
plane centering are invariant under simultaneous spin flip.

It remains to estimate all terms that can see the wall, including
saved scalar increments.  Every such source footprint contains the
mark and reaches the boundary region.  At scale $j$ its carrier length
is therefore at least $c n/r_j$.  The deterministic weighted source
transport costs a fixed factor $K$ per step, even with arbitrary
principal background sizes: selected averages are normalized, retained
large terms do not enlarge a source norm, and the small-field
coefficient sums are bounded.  This is the source estimate of
Section~\ref{sec:norms} applied without averaging over the randomness.
Keep an additional exponential carrier weight.  At each possible
scale of first boundary contact it supplies
$\exp(-c'n/r_j)$.  Sum these contributions, retaining the source
rescaling up to $J_n$; changing $K$ absorbs both that rescaling and
the sum over the possible first-contact steps.  The total deterministic
cost at the stopping scale is at most
\begin{equation}
 C K^{J_n}\exp(-c'n/r_{J_n})
 \le \exp\bigl(C\log n-c''(\log n)^2\bigr).
 \label{read:boundary-tail}
\end{equation}
One can equivalently subtract the plane and finite-box constructions:
the first coefficient where they differ must have a footprint meeting
the wall, and each later difference is transported with the same
weighted estimate.  This verifies the bound for all nonbulk
coefficients without assuming that their laws agree with plane laws.
For an odd scalar the carrier weight is used before its support is
forgotten, so \eqref{read:boundary-tail} includes all nonzero scalar
contributions.  Initial fixed constants merely change $C$.

The exact normalized one-source identity now expresses the plus
magnetization, multiplied by $r_{J_n}^d$, as the terminal expectation
of the scaled source plus its saved scalar.  Its absolute value is
bounded by the total sup costs just estimated, irrespective of the
terminal background.  Hence, for some fixed $M<\infty$,
\[
 \norm{\langle\sigma_y\rangle^+_{n,*,\omega}}_{L^{400}}
 \le r_{J_n}^{-d}\left(C(1+J_n)^M
                  +C K^{J_n}e^{-c'n/r_{J_n}}\right)
       =n^{-d+o(1)}.
\]
Raising this inequality to the power $400$ proves
\eqref{read:arm-bound}.  Translating the box changes none of the
constants because the source estimates are uniform in its mark.
Moving the wall by a bounded number of lattice layers changes the
lower bound on the boundary-reaching carrier length only by a fixed
factor, proving the remaining conventions.
\end{proof}

In particular the averaged plus magnetization vanishes at tuning.
The next section uses the stronger finite-box bound
\eqref{read:arm-bound} to identify the tuned parameter and to show
that magnetization becomes positive above it.
